\chapter{Deployment View}

\section{Resource Limits}
No resource limits are identified at the scope of the current Architecture design.
\\
\\
The SUM\_TCS is developed in order to be deployed in a 32-bit automotive microcontroller-based system (Single Core).

\section{AUTOSAR Resources}
The current architecture has been designed in order to optimize the memory usage and the timing resources according to AR safety aspects and general specification.
\\
\\SUM\_TCS architecture has been developed in order to be fully compliant with Classic AR.
\\
\\
All AUTOSAR resources that shall be deployed for the implementation and the integration of SUM\_TCS are presented in the previous sections of the current document.

\section{Bus System}
SUM\_TCS is originally designed in order to communicate in a multi-channel CAN and CAN-FD Bus System; Although The Design of SUM\_TCS is totally transparent to 
the Bus system configuration and implementation.

\section{Interrupts}
N/A

\section{Watchdog}
N/A

\section{External hardware dependencies}
N/A

\section{OS}
The SUM\_TCS is designed to interfact into a RTOS.
\\
\\
The SUM\_TCS Runnables shall be mapped in OS Tasks to be properly scheduled by integrator.
\\
OS is assumed to be configured as \textbf{non-preemptive}.
\\
\\
SUM\_TCS defines:
\begin{itemize}
	\item 2 \textit{Periodic} Runnables;
	\item 1 \textit{Data-Received} triggered Runnable;
	\item 3 \textit{Operation-Invoked} triggered Runnable;
\end{itemize}

It is strongly suggested to map the periodic runnable in a Task and the Data-Received triggered one in another.
\\
\\
It is strongly suggested to give an higher priority to Task for Data-Received event than Periodic Task.
\\
\\
Within the Periodic task, it is strongly suggested to give an higher priority to \hyperref[subsec:Runnable_TCS_MasterClockPolling]{Runnable_TCS_MasterClockPolling} 
than \hyperref[subsec:Runnable_TCS_TimeManagement]{Runnable_TCS_TimeManagement}.

